\appendix
\section{Exploratory Quantum Experiment}
\label{app:quantum}
We use IBM Qiskit for a synthetic-only experiment outside the clinical runtime.
We balance supported claims against unsupported and contradicted claims, using
\QuantumTrain{} training and \QuantumTest{} test examples from frozen verifier
splits. The binary task differs from the main verifier's label space. Pair features
combine cosine similarity, products and absolute differences of character-hash
vectors. We fit PCA and scaling on training features and encode
\QuantumQubits{} components in a ZZ feature map with \QuantumReps{} repetitions.
The intended Gemini feature experiment was blocked by quota, so we report the
offline character-hash control and do not infer semantic-embedding performance.

On \texttt{ibm\_marrakesh}, we execute \QuantumCircuits{} upper-triangular overlap
circuits, including diagonals, for a \QuantumBlockSize{} by \QuantumBlockSize{}
test-kernel block. We use \QuantumShots{} shots per circuit. Median transpiled depth
is \QuantumDepth{} and recorded execution time is \QuantumSeconds{} s. The mean
absolute error relative to the exact simulator block is \QuantumMae{}. These
measurements describe kernel fidelity, not classification accuracy. A test-to-test
block cannot replace the classifier's test-to-training kernel, so we make no
held-out hardware classifier claim. The full aggregate and backend provenance remain
in \texttt{ml/quantum/results.json}.

We also implement QAOA evidence selection on small synthetic candidate sets, compared
with exact and greedy search. This is a relevance surrogate rather than a clinical
standard, and exact enumeration is practical at this scale. The experiments establish
neither scaling behaviour, quantum advantage nor transfer to patient care.
